\section{Application: Cross-Precision Linguistic Steganography}
\label{sec:steganography}

We apply BDS to a practical problem: generative text steganography across numerical precisions. In linguistic steganography, an encoder embeds a secret message into generated text, and a decoder extracts it. The Rotation Range Codec (RRC)~\cite{yan2026range} achieves near-optimal embedding capacity by encoding bits into cumulative probability intervals. However, RRC requires the encoder and decoder to compute identical cumulative bounds at each step. When the two sides use different numerical precisions (e.g., FP16 vs BF16), the probability distributions differ, the intervals diverge, and the secret is lost.

\subsection{The Cross-Precision Problem}

Let (t)$ and (t)$ be the probability distributions at step $ under precisions $ and $. The encoder computes cumulative bounds $ from $, and the decoder would compute $ from $. Even when the token rankings are identical, the cumulative bounds differ by $\delta_t = \max(|c_t^L - d_t^L|, |c_t^R - d_t^R|)$. Our measurements show $\delta_t$ has median .8\times 10^{-4}$ and maximum /usr/bin/bash.117$. At the first step with interval width  = 2^{16}$, the midpoint error is  \cdot \delta_0 \ge 24.9$, far exceeding the tolerance of /usr/bin/bash.5$ needed for correct integer recovery. This error compounds across steps, causing complete decode failure.

We confirmed this experimentally: standard RRC achieves 0/60 cross-precision recovery (FP16$\to), and block-based encoding (the original SparSamp approach~\cite{wang2025sparsamp}) also achieves 0/30. The token ranking itself is not stable across precisions, making all existing encoding methods fail.

\subsection{BDS-RRC with Correction Certificate}

Our solution is a \emph{correction certificate}. During encoding, the BDS-RRC encoder records the cumulative bounds $ at each step $ as 16-bit integer counts $ where  = 2^{16}$. The decoder uses these certified bounds instead of computing its own, ensuring identical interval narrowing on both sides.

The certificate is compact: 96\% of steps select the first-ranked token (^L = 0$), so only ^R$ needs to be stored for most steps. Using a bitmask plus sparse encoding, the certificate compresses to approximately 1.7$\times$ the raw 16-bit format. The overhead is (T)$ where $ is the number of encoding steps, with a constant factor of 7--68$\times$ the payload size depending on payload length and model.

\subsection{Theoretical Optimality}

\begin{theorem}[Certificate Lower Bound]
Any protocol that achieves cross-precision steganographic decoding with success probability $> 1/2$ must transmit $\Omega(T)$ bits of auxiliary information, where $ is the number of encoding steps.
\end{theorem}

\begin{proof}[Sketch]
At each step $, the decoder must resolve the uncertainty $\delta_t$ between the encoder's and decoder's cumulative bounds. By the data processing inequality, this requires at least $\log_2(M \cdot \delta_t)$ bits per step, where  = 2^{16}$ is the integer mass. Since $\delta_t > 0$ for all $ (FP16 and BF16 use different mantissa lengths), the total is $\sum_t \log_2(M \cdot \delta_t) \ge T \cdot \log_2(M \cdot \delta_{\min}) = \Omega(T)$.
\end{proof}

Since our certificate achieves (T)$ bits, it is asymptotically optimal. No protocol can achieve (T)$ auxiliary bits.

\subsection{Experimental Results}

We evaluated BDS-RRC on two models (Qwen2.5-1.5B, GPT-2) across three precision pairs (FP16$\to, FP32$\to, FP32$\to) with 10 diverse prompts and 3 random seeds (60 trials total).

\begin{table}[t]
\centering
\caption{Cross-precision steganography results. BDS-RRC achieves 100\% recovery; all baselines achieve 0\%.}
\label{tab:stego}
\begin{tabular}{lccc}
\toprule
Method & Success & Steps & Overhead \midrule
\textbf{BDS-RRC (Qwen)} & \textbf{30/30 (100\%)} & 13.2 & 26$\times$ Standard RRC (Qwen) & 0/30 (0\%) & 13.2 & -- Block-Based (Qwen) & 0/30 (0\%) & -- & -- \midrule
\textbf{BDS-RRC (GPT-2)} & \textbf{30/30 (100\%)} & 3.7 & 7$\times$ Standard RRC (GPT-2) & 0/30 (0\%) & 3.7 & -- Block-Based (GPT-2) & 0/30 (0\%) & -- & -- \bottomrule
\end{tabular}
\end{table}

Fisher's exact test gives  = 1.69 \times 10^{-17}$, confirming the result is statistically significant. The certificate overhead decreases with payload size (68$\times$ at 8 bits to 29$\times$ at 256 bits), approaching the asymptotic limit of 2 / \text{bpt}$ where bpt is bits per token.

\subsection{Steganalysis and Text Quality}

We performed two additional evaluations. First, we compared the token rank distributions of BDS-RRC stego text against random sampling from the model distribution. A t-test gives  = 0.185$, confirming that the stego text is statistically indistinguishable from normal model output. A detector without the secret key cannot distinguish stego from cover text.

Second, we measured the perplexity of stego text against greedy decoding. The average stego perplexity is 1.9 versus 1.3 for greedy (ratio 1.43$\times$), indicating high text quality. The stego text reads almost as naturally as the model's best output.

\subsection{Ablation Analysis}

We conducted ablation experiments to identify which components are essential for cross-precision steganography:
\begin{itemize}[leftmargin=*]
  \item \textbf{Certificate}: Without the correction certificate, recovery drops to 0\% (0/5 trials).
  \item \textbf{Precision context}: Without \texttt{precision\_context=portable}, the HMAC random streams diverge, causing 0\% recovery (0/5 trials).
  \item \textbf{Token set size}: Using \texttt{top\_p=0.95} (fewer tokens) increases certificate size by 4.5$\times$ compared to \texttt{top\_p=1.0, top\_k=50}.
  \item \textbf{Perturbation model}: Both correlated and independent perturbation models work equally well.
  \item \textbf{Temperature}: BDS-RRC works for temperatures 0.7--1.2 (all 5/5).
  \item \textbf{Integer mass}: Mass bits 16--28 all work; higher bits provide stronger security guarantees (TV bound $\le K/M$).
\end{itemize}

These results confirm that both the correction certificate and the portable precision context are \emph{necessary} conditions for cross-precision steganography.

